\documentclass[11pt]{article}
\usepackage{preamble}
\renewcommand{\answer}[1]{}

\begin{document}

\ladderheading{AP Statistics \textperiodcentered\ Unit 1 \textperiodcentered\ Topic 1.8 \textperiodcentered\ B: 2026-09-22 \textperiodcentered\ E: 2026-09-28}{Two Displays, One Data Set}

\focusbanner{TODAY'S PROBLEM}

The displays show practice-test times, in minutes, for 24 students. Two students interpret the boxplot.\par\smallskip \textbf{Omar:} ``Most students took about 32 minutes---the middle of the box.''\par \textbf{Nia:} ``A boxplot cannot tell us whether most students took about 32 minutes.''\par
\begin{center}
\begin{tikzpicture}
\begin{axis}[
  width=0.92\linewidth,
  height=1.85in,
  xmin=20, xmax=58,
  ymin=-1.6, ymax=3.8,
  ytick=\empty,
  axis y line=none,
  axis x line=bottom,
  xtick={22,24,32,40.5,55},
  xticklabels={22,24,32,40.5,55},
  tick label style={font=\scriptsize},
  xlabel={Practice-test time (minutes)},
  label style={font=\small},
  clip=false
]
\node[anchor=west,font=\scriptsize\bfseries] at (axis cs:20.4,3.35) {Dotplot};
\addplot[only marks, mark=*, mark size=2.4pt, color=dayblue] coordinates {
  (22,1) (22,2) (23,1) (23,2) (24,1) (24,2) (24,3)
  (25,1) (25,2) (25,3) (26,1) (26,2)
  (38,1) (38,2) (39,1) (39,2) (40,1) (40,2)
  (41,1) (41,2) (42,1) (43,1) (44,1) (55,1)
};
\node[anchor=west,font=\scriptsize\bfseries] at (axis cs:20.4,-0.15) {Boxplot};
\draw[dayblue,thick] (axis cs:22,-0.85) -- (axis cs:55,-0.85);
\draw[dayblue,thick] (axis cs:22,-1.15) -- (axis cs:22,-0.55);
\draw[dayblue,thick] (axis cs:55,-1.15) -- (axis cs:55,-0.55);
\filldraw[fill=calloutblue,draw=dayblue,thick]
  (axis cs:24,-1.25) rectangle (axis cs:40.5,-0.45);
\draw[dayblue,very thick] (axis cs:32,-1.25) -- (axis cs:32,-0.45);
\end{axis}
\end{tikzpicture}
\end{center}
\begin{firsttakebox}{FIRST TAKE --- before the video (5 min)}
Who do you trust more, Omar or Nia? Point to one feature in either display.
\workspace{0.9in}
\end{firsttakebox}

\begin{callout}{\IconBook\ DOTPLOTS AND BOXPLOTS (keep on the page)}{calloutgreen}
A \textbf{dotplot} keeps every observation, so it can show clusters, gaps, peaks, and isolated values.\
A \textbf{boxplot} compresses the data to the five-number summary: minimum, Q1, median, Q3, maximum
(with separately marked outliers). The box contains the middle $50\%$ of observations.\
A boxplot shows center and spread efficiently, but it may hide modality, clusters, and gaps.
\end{callout}

\newpage
\textbf{Finish after the video:}\par\smallskip

\dokbadge{1}~\textbf{(a)} Read the median and the first and third quartiles from the boxplot. What percent of the times lie in the box from Q1 to Q3?\par
\workspace{0.7in}

\dokbadge{2}~\textbf{(b)} Describe the dotplot in context: include the two clusters, the gap, and the isolated high value. State two features that the boxplot preserves.\par
\workspace{1.4in}

\dokbadge{3}~\textbf{(c)} Whose interpretation is better supported? Cite the dotplot feature that settles it. Explain what a boxplot can and cannot show here. A coach asks whether most students need more than 30 minutes: choose a display, and explain what the boxplot alone would lead the coach to believe.\par
\workspace{2.0in}

\begin{sentenceframebox}\raggedright\textbf{Frame:} \blank[0.8in] is better supported because the dotplot shows \blank[2.2in] around the median of 32 minutes.\end{sentenceframebox}

\begin{sentenceframebox}\raggedright\textbf{Frame:} The boxplot shows \blank[1.4in], but it hides \blank[1.8in], so the coach should use the \blank[0.8in].\end{sentenceframebox}

\begin{summaryexitbox}{EXIT --- turn this in}
One thing the video changed about my first take: \hrulefill\par
\workspace{0.4in}
\end{summaryexitbox}

\end{document}
